\section{相关工作}
\subsection{分子表示与条件溶解度建模}
分子描述符通过分子量、脂溶性、极性表面积和氢键能力等固定统计量表征结构；扩展连接指纹则通过局部原子邻域编码分子连接模式\cite{ecfp2010}。分子预训练进一步从大量无标签结构学习表示，例如MolCLR利用对比学习改善分子表征\cite{molclr2022}。对于多溶剂预测，单独表达溶质并不充分，MolMerger通过显式溶质--溶剂相互作用构造图表示\cite{molmerger2024}，说明两个分子之间的关系本身具有建模价值。

温度条件提供了另一类信息。FASTSOLV将分子表示、溶剂和温度联合建模，并利用温度梯度监督改善响应学习\cite{fastsolv2025}。本文据此以相邻温度观测差分补充逐点监督，并将其与母体均衡、结构交互和显式温度响应结合。

溶解度数据研究同时强调测量条件、离子状态和数据重叠对模型评价的影响\cite{solubility2024,bigsoldb2025}。这使分组划分和条件信息处理成为模型比较的重要组成部分。

近年来，分子表示学习进一步向几何信息和多模态知识扩展。Wang等总结了自监督学习与几何深度学习在科学发现中的作用，强调结构信息与数据表示之间的联系\cite{discovery2023}。MoleculeSTM通过对比学习对齐分子结构和文本描述，支持基于文本的分子检索与编辑\cite{moleculestm2023}，为结构分支引入化学语义提供了进一步方向。

\subsection{小样本表格学习与分组评价}
岭回归通过二次正则化控制参数方差\cite{ridge1970}；随机森林和ExtraTrees利用树集成拟合非线性关系\cite{randomforest2001,extratrees2006}；梯度提升通过逐步拟合残差优化目标函数\cite{boosting2001}。这些方法对中低维描述符具有直接适用性，也构成检验新模型的必要基线。

深度表格学习的近期进展表明，训练配置和评价协议会影响树模型与神经网络之间的比较\cite{tabular2021,trees2022}。RealMLP结合数值预处理、参数化和优化策略构建竞争力较强的MLP\cite{realmlp2024}；TabM通过参数高效集成产生多个成员预测\cite{tabm2025}。TabPFN进一步展示了表格基础模型在小数据任务上的潜力\cite{tabpfn2025}。本文选用TabM与RealMLP作为近期基线，并设置同规模MLP对照，以比较架构、条件监督和母体均衡的作用。

当同一对象在多个条件下重复出现时，随机划分记录可能使训练与验证共享对象信息。科学机器学习中的数据泄漏研究指出，评价单位应与实际推广任务匹配\cite{leakage2023}。因此，本文将药物母体作为划分单位，使用母体等权误差作为主要评价指标，并将训练损失中的均衡原则与之对应。

\subsection{模型解释与因果敏感性}
SHAP将预测分解为特征贡献，TreeSHAP使该思想能够高效应用于树模型\cite{shap2017,treeshap2020}；积分梯度则沿基准输入到观测输入的路径积分梯度，为神经网络提供输入归因\cite{integratedgradients2017}。本文分别解释混合模型的结构树分支与交互网络，并以整组置换比较完整预测函数对药物、溶剂和温度的依赖。

预测归因回答模型如何利用输入，因果分析还涉及干预定义与混杂控制\cite{causal2021}。本文以组合固定效应研究温度关联，再使用偏$R^2$参数化的遗漏变量偏差分析\cite{sensitivity2020}，定量说明未观测因素需要多强的关联才能改变该系数。
